\documentclass[11pt]{article}

% --- Series style ---
\usepackage{series}

% --- Math and layout ---
\usepackage{amsmath,amssymb,amsthm}
\usepackage{mathtools}
\usepackage{geometry}
\usepackage{hyperref}

\geometry{margin=1in}


% --- Metadata ---
\title{Photons, Entanglement, and Null Updating in Modal Triplet Theory}
\author{Peter Nero}
\date{January 2026}

\begin{document}
\maketitle

\begin{abstract}
We present a unified structural account of photons in Modal Triplet Theory (MTT). In MTT,
physical description arises from admissible coherent configurations on a higher-dimensional
modal space, projected to an effective four-dimensional encoding. Within this framework,
photons are identified as massless gauge-coherence modes whose preferred physical encoding
is global and non-factorizing. Localization, particle-like behavior, and discrete detection
events emerge only when admissibility constraints enforce a partition of this encoding, as
in interaction or measurement.

We show that null propagation at the speed of light arises as the unique admissible updating
of massless gauge coherence under stability and contractivity conditions, rather than as a
kinematic postulate. Wave--particle duality is reinterpreted as a choice of encoding valid on
different admissible domains. Gravitational lensing and redshift are described as continuous
re-encodings of the same stretched coherence constraint under varying admissibility geometry,
without invoking forces acting on localized particles.

Photons near horizons provide a particularly transparent illustration of the framework.
Horizons are identified as admissibility barriers across which the observable map loses a
global right inverse, forcing exterior re-encoding. The resulting thermal description is
shown to be consistent with Hawking's equations and is interpreted as a basin-measure shadow
of restricted encoding, structurally analogous to measurement-induced probabilities. Time
is treated as an encoding-dependent construct associated with coherence rebalancing and
record formation; massless gauge modes admit atemporal encodings along null updating
directions.

This paper does not modify the equations of quantum electrodynamics, general relativity, or
quantum field theory on curved spacetime. Rather, it clarifies their structural origin within
MTT and unifies photon propagation, entanglement, measurement, horizons, and thermality as
manifestations of a single admissibility--encoding framework.
\end{abstract}

% ============================================================
%  What Is a Photon?
%  Coherent Encoding and Null Propagation in Modal Triplet Theory
% ============================================================

\section{Introduction}
\label{sec:introduction}

Photons occupy a central role in quantum theory, general relativity, and quantum field
theory, yet their ontological status remains conceptually fragmented. In standard quantum
mechanics, photons are treated as excitations of the electromagnetic field whose states may
be delocalized, entangled, or sharply localized depending on experimental context. In
relativistic field theory, they are massless gauge bosons propagating along null directions.
In gravitational settings, photons undergo lensing, redshift, and horizon-induced thermality.
Despite the mathematical consistency of these descriptions, they are rarely unified into a
single structural picture.

Modal Triplet Theory (MTT) provides a framework in which these disparate descriptions arise
as different encodings of a single underlying coherent structure. In MTT, the fundamental
object is not a particle propagating through spacetime, but a coherent modal configuration
defined on a higher-dimensional configuration space and projected into an effective
four-dimensional description. Within this framework, photons emerge as massless
gauge-coherence modes whose natural encoding is global and non-factorizing.

The purpose of this paper is to give a precise account of what a photon \emph{is} in MTT, and
to show how null propagation, wave--particle duality, entanglement, gravitational lensing,
redshift, horizons, and Hawking radiation arise as consequences of admissibility and encoding,
rather than as independent postulates. The treatment is not intended to modify the equations
of quantum electrodynamics or general relativity, but to clarify their structural origin and
interpretation within the MTT framework.

In particular, we argue that:
\begin{itemize}
\item The preferred encoding of a photon is a globally stretched, entangled gauge-coherence
state.
\item Null propagation arises as the only admissible updating of such coherence under the
fundamental contractivity and stability conditions.
\item Localization and ``particle-like'' behavior appear only when measurement or interaction
enforces a partition of the coherent encoding.
\item Gravitational effects on photons correspond to continuous re-encoding of the same
coherence constraint under varying admissibility geometry.
\item Horizons represent admissibility barriers where a global encoding fails, yielding
thermal exterior descriptions consistent with Hawking radiation.
\end{itemize}

This paper is intended as a structural synthesis and interpretive clarification, building on
the technical results of the MTT$\Rightarrow$QM, MTT$\Rightarrow$GR, measurement, entanglement,
and horizons papers, and situating photons within the broader admissibility--encoding
framework.

---

\section{Coherent Encoding of Gauge Fields}
\label{sec:coherent_gauge_encoding}

\subsection{Modal configuration space and projection}
\label{sec:modal_space_projection}

In Modal Triplet Theory, physical description begins with a modal configuration space
\begin{equation}
M_{10} = Y^{4} \times B_1 \times B_2 \times B_3 ,
\end{equation}
where $Y^{4}$ is the emergent four-dimensional spacetime and $B_i$ are compact internal
modal fibers. Physical states are not identified with arbitrary configurations on $M_{10}$,
but with those lying in the coherent sector defined by the joint harmonic projector
$\Pi = \Pi_{B_1}\Pi_{B_2}\Pi_{B_3}$.

The observable map
\begin{equation}
P := I \circ \Pi
\end{equation}
projects coherent modal configurations to effective four-dimensional states. This map is
many-to-one: distinct modal configurations may project to the same observable state. As a
result, the effective description is intrinsically compressed, and different encodings of
the same physical situation may coexist on different admissible domains.

\subsection{Photons as massless gauge-coherence modes}
\label{sec:photon_definition}

Within this framework, a photon corresponds to a massless gauge-coherence mode of the
coherent sector. Such a mode has two defining properties:

\begin{enumerate}
\item It carries no internal coherence inertia: there is no mass term or internal
rebalancing cost associated with maintaining the mode.
\item Its admissibility is governed entirely by gauge coherence constraints and the
stability of the coherent projector.
\end{enumerate}

Because of the absence of internal inertia, there is no preferred localization scale for a
photon in the coherent sector. The natural encoding of a photon is therefore not a localized
particle state but a globally stretched coherence constraint whose support may extend across
large null-connected regions of spacetime.

In this sense, a single-photon state is generically entangled across spacetime modes
(direction, frequency, polarization), and factorized descriptions correspond to special
encodings imposed by additional constraints.

\subsection{Null propagation as admissible updating}
\label{sec:null_propagation}

Although the preferred encoding of a photon is global, coherence constraints must still be
updated consistently with the modal dynamics. Admissibility imposes strong restrictions on
how this updating may occur.

Static gauge-coherence configurations are generically inadmissible: without updating, they
violate stability conditions and erode the contractivity margins required for a coherent
sector to persist. Superluminal updating is likewise inadmissible, as it would violate the
locality and boundedness assumptions underlying the modal dynamics.

The only admissible updating for a massless gauge-coherence mode is therefore along null
directions of the emergent spacetime geometry. Null propagation is not imposed as a kinematic
postulate but arises as the unique way to update a globally stretched coherence constraint
without destroying admissibility.

Thus photon propagation at speed $c$ is reinterpreted as:
\begin{quote}
\emph{Null updating of a globally entangled gauge-coherence encoding under admissibility
constraints.}
\end{quote}

This interpretation reproduces all standard predictions of Maxwell theory and quantum
electrodynamics while clarifying why localization and propagation are conceptually distinct.

\subsection{Position, momentum, and conditional localization}
\label{sec:position_momentum}

In this picture, momentum labels the orientation and phase structure of the null updating,
and is intrinsic to the stretched coherence encoding. Position, by contrast, is not an
intrinsic property of the photon prior to interaction or measurement. Localization requires
an enforced partition of the encoding, typically associated with record formation or strong
interaction with matter.

Accordingly, wave--particle duality is understood as an encoding choice:
\begin{itemize}
\item Wave-like behavior corresponds to the globally stretched, non-factorizing encoding.
\item Particle-like behavior corresponds to a localized encoding forced by admissibility
constraints near measurement or interaction boundaries.
\end{itemize}

No contradiction arises, as these encodings are valid on different admissible domains and
cannot generally be maintained simultaneously.
% ============================================================
%  Supplementary Photon Subsections:
%  Polarization/Helicity, Virtual Photons, and Causality
% ============================================================

\subsection{Polarization and Helicity}
\label{sec:polarization_helicity}

Photons carry spin-$1$ and are observed to possess two physical polarization states,
corresponding to helicity $\pm 1$. In the Modal Triplet Theory framework, polarization and
helicity are not properties of a localized particle but internal orientations of the
gauge-coherence encoding under admissible null updating.

Helicity conservation reflects the topological stability of massless gauge coherence:
for a null-updating mode with no internal rebalancing cost, only two transverse orientations
are admissible. Parallel transport of polarization is therefore governed by admissibility
geometry rather than by forces acting on a localized spin degree of freedom. Effects such as
gravitational Faraday rotation and polarization transport in curved spacetime correspond to
continuous re-orientation of the internal gauge-coherence structure under admissible null
updating.

This interpretation preserves all standard results regarding photon polarization while
clarifying that helicity is a global encoding constraint rather than an intrinsic property
of a point-like object.

\subsection{Virtual Photons and Perturbative Encodings}
\label{sec:virtual_photons}

In perturbative quantum field theory, interactions are often represented in terms of
``virtual photons.'' Within the MTT framework, such objects do not correspond to admissible
coherent gauge modes. Instead, they are internal elements of a \emph{perturbative encoding}
used to approximate interaction constraints within a particular calculational scheme.

Virtual photons do not propagate as physical degrees of freedom, do not obey null updating,
and are not subject to admissibility requirements in the same sense as real photons. They
are bookkeeping devices that arise when interactions are expanded around a factorized
encoding of the coherent sector. Consequently, no ontological significance is assigned to
virtual photons beyond their role in organizing perturbative corrections.

This distinction avoids common interpretive confusions and reinforces the separation
between physical coherent encodings and auxiliary calculational constructs.

\subsection{Locality, Causality, and Global Encoding}
\label{sec:causality_locality}

The use of globally stretched and entangled encodings does not compromise locality or
causality in the effective four-dimensional description. Admissibility enforces that all
coherence updates occur locally and propagate causally with respect to the emergent
spacetime geometry.

In particular:
\begin{itemize}
\item Operator commutators vanish outside the light cone in the effective description,
preserving microcausality.
\item No superluminal signaling is possible, even though the preferred encoding may be
globally non-factorizing.
\item Apparent instantaneous correlations reflect global consistency of a single coherent
configuration rather than physical influence.
\end{itemize}

Thus global encoding should not be conflated with global interaction. The distinction
between encoding scope and dynamical updating is essential: coherence may be shared
globally, but its admissible evolution remains strictly local.


% ============================================================
%  Section 6: Gravitational Lensing and Redshift
% ============================================================

\section{Gravitational Lensing and Redshift as Re-encoding}
\label{sec:lensing_redshift}

\subsection{Spacetime geometry as admissibility bookkeeping}
\label{sec:geometry_bookkeeping}

In Modal Triplet Theory, spacetime geometry does not enter as a fundamental interaction
acting on particles, but as an emergent bookkeeping structure encoding which coherence
updates are admissible. Curvature, light cones, and geodesics summarize how coherent
configurations may be transported without violating stability, boundedness, and the
fundamental contractivity condition.

In particular, null directions in the emergent metric encode the admissible updating
directions for massless gauge-coherence modes. Variations in spacetime geometry therefore
correspond to variations in admissibility structure, rather than to forces acting on
localized objects.

\subsection{Gravitational lensing}
\label{sec:lensing}

Gravitational lensing is commonly described as the deflection of photon trajectories by a
gravitational field. In the MTT framework, this description is replaced by a purely
structural statement:

\begin{quote}
\emph{Gravitational lensing is the continuous deformation of admissible null updating
directions for a globally stretched gauge-coherence encoding.}
\end{quote}

Because a photon is encoded as a massless gauge-coherence constraint with global support,
it does not ``experience'' local forces or impulses. Instead, the null directions along
which coherence may be updated bend in response to the local admissibility geometry. The
photon's stretched encoding must follow these admissible null directions, yielding the
observed lensing effect.

This interpretation immediately explains several robust features of gravitational
lensing:
\begin{itemize}
\item Lensing is achromatic, as it depends only on admissible null structure and not on
internal photon properties.
\item Phase coherence and interference can be preserved over cosmological distances,
since the encoding remains globally coherent.
\item Polarization transport follows geometric rules, reflecting the transport of gauge
coherence rather than particle deflection.
\end{itemize}

No additional interaction mechanism is required beyond the admissibility constraints that
define the emergent spacetime geometry.

\subsection{Gravitational redshift}
\label{sec:gravitational_redshift}

Gravitational redshift is traditionally described as a photon losing energy as it climbs
out of a gravitational potential. In the MTT framework, this language is replaced by a
relative-encoding interpretation.

Frequency is not an intrinsic property of a massless gauge-coherence mode. Instead, it is
defined relationally, with respect to local clocks. Clocks, in turn, are massive systems
whose internal coherence requires time-like rebalancing and therefore provides a local
time scale.

When a photon is emitted and later absorbed in regions with different admissibility
geometry, the same globally stretched coherence constraint is re-encoded relative to
different local clock structures. The mismatch between these encodings appears as a
frequency shift.

Thus:
\begin{quote}
\emph{Gravitational redshift arises because the same coherent photon encoding is referenced
to different local time encodings at emission and detection.}
\end{quote}

No intrinsic change occurs in the photon itself. The redshift reflects a change in the
encoding context, not an energy loss along the photon's path.

\subsection{Cosmological redshift}
\label{sec:cosmological_redshift}

Cosmological redshift admits an analogous interpretation. As the admissibility geometry
evolves with cosmic expansion, the null updating structure of spacetime is continuously
re-parameterized. A photon propagating across cosmological distances remains a single
stretched coherent encoding, but its frequency is re-expressed relative to the evolving
clock structure of comoving observers.

Cosmological redshift is therefore a cumulative re-encoding effect associated with the
global evolution of admissibility geometry, rather than a gradual stretching of a
localized wave or loss of photon energy.

\subsection{Persistence of coherence over large distances}
\label{sec:coherence_persistence}

A natural question is why photon coherence can persist over extremely large distances,
even in the presence of curvature and expansion. The MTT answer is straightforward: massless
gauge-coherence modes carry no internal rebalancing cost and therefore accumulate no
intrinsic decoherence during admissible propagation.

Decoherence arises only through interactions that force partitioning of the encoding
(scattering, absorption, or measurement). In the absence of such interactions, the
globally stretched encoding remains admissible and coherent regardless of distance or
elapsed coordinate time.

\subsection{Summary}
\label{sec:lensing_redshift_summary}

Gravitational lensing and redshift are thus unified within a single MTT mechanism:
\begin{itemize}
\item Lensing reflects the deformation of admissible null updating directions.
\item Redshift reflects the re-encoding of the same global coherence relative to different
local clock structures.
\item Neither effect requires intrinsic changes to the photon or particle-like forces.
\end{itemize}

Both phenomena follow inevitably from the interpretation of spacetime geometry as
admissibility bookkeeping for coherent encodings.
% ============================================================
%  Section 7: Photons Near Horizons and Hawking Radiation
% ============================================================

\section{Photons Near Horizons and Hawking Radiation}
\label{sec:horizons_hawking}

\subsection{Horizons as admissibility barriers}
\label{sec:horizons_admissibility}

In the MTT framework, a black hole horizon is not primarily a locus of divergent curvature
or violent local physics. Instead, it is an \emph{admissibility barrier}: a codimension-one
boundary across which the observable map
\begin{equation}
P = I \circ \Pi
\end{equation}
ceases to admit a global right inverse on the full algebra of coherent configurations.

Outside the horizon, a coherent encoding exists that supports a consistent four-dimensional
description. Inside the horizon, coherent configurations continue to exist at the modal
level, but they are no longer jointly encodable with the exterior in a single admissible
chart. This loss of global encodability is the defining structural feature of a horizon in
MTT.

\subsection{Photon coherence approaching a horizon}
\label{sec:photon_near_horizon}

As discussed in previous sections, a photon is naturally encoded as a globally stretched,
massless gauge-coherence constraint whose admissible updating occurs along null directions.
Far from a horizon, this encoding remains stable and coherent across large spacetime regions.

As the photon approaches a horizon, the admissibility geometry changes. While local null
updating remains well-defined, the global encoding that spans both the exterior and interior
regions becomes marginal. In MTT terms, the stretched photon encoding enters a
\emph{selection-front regime}, where stability margins shrink and the projector loses
regularity across the horizon.

Crucially, nothing singular occurs in the underlying modal dynamics. The breakdown is not in
the evolution of coherence itself, but in the ability of the exterior encoding to represent
that coherence globally.

\subsection{Infinite redshift as encoding mismatch}
\label{sec:infinite_redshift}

From the perspective of an exterior observer, photons approaching a horizon exhibit an
infinite redshift. In the MTT framework, this effect has a clear encoding interpretation.

Frequency is defined relative to local clocks, which are massive systems with internal
coherence requiring time-like rebalancing. Near a horizon, the admissibility geometry causes
the relation between exterior clock encodings and the stretched photon coherence to become
singular. The same global coherence constraint is therefore re-encoded as arbitrarily low
frequency when referenced to exterior time.

This infinite redshift reflects the failure of a single encoding to remain valid across the
horizon, rather than any intrinsic change in the photon.

\subsection{Forced re-encoding and the emergence of thermality}
\label{sec:forced_reencoding}

Once the stretched photon encoding crosses the admissibility barrier associated with the
horizon, the exterior description must adopt a restricted encoding that excludes interior
degrees of freedom. The exterior observer's algebra is therefore incomplete, and the global
coherent configuration must be re-expressed in terms of exterior-admissible basins.

This situation is structurally identical to measurement-induced collapse: in both cases, a
restriction of admissibility forces a re-encoding of a globally coherent configuration into
a set of mutually exclusive, record-compatible outcomes.

In the horizon case, the admissible exterior basins are labeled continuously by asymptotic
frequency $\omega$. The weights assigned to these basins are given by the same universal
basin-measure functional that governs measurement probabilities. In the stationary,
semiclassical regime, this yields
\begin{equation}
p(\omega) \propto e^{-\beta \omega},
\end{equation}
with $\beta$ identified with the inverse Hawking temperature.

\subsection{Consistency with Hawking's equations}
\label{sec:hawking_consistency}

The above interpretation does not modify the standard equations used to derive Hawking
radiation. The wave equation on a fixed black hole background, the Bogoliubov transformation
relating in- and out-modes, and the resulting thermal spectrum remain unchanged.

In MTT terms, these equations describe the behavior of the exterior encoding near an
admissibility barrier. The appearance of a thermal density matrix is a direct consequence of
restricting a globally coherent configuration to a subalgebra that does not admit a global
inverse under $P$.

Thus Hawking radiation is understood as the horizon-barrier instance of a general MTT
principle: whenever admissibility forces a non-invertible projection, the effective
description becomes probabilistic, with weights determined by the invariant basin measure.

\subsection{Information, unitarity, and irreversibility}
\label{sec:information_unitarity}

At the modal level, evolution remains deterministic and unitary. No information is destroyed
in the full configuration space. The apparent loss of information associated with Hawking
radiation arises entirely from the use of a restricted encoding that cannot represent all
degrees of freedom simultaneously.

This is the same structural mechanism responsible for the appearance of irreversibility and
probability in measurement processes. Horizons and measurements are therefore unified as
instances of admissibility barriers, differing only in the nature of the restriction
(geometric versus instrumental).

\subsection{Summary}
\label{sec:horizons_summary}

Photons near horizons provide a particularly clear illustration of the MTT encoding
framework:
\begin{itemize}
\item Photons remain globally coherent massless gauge modes at the modal level.
\item Horizons correspond to admissibility barriers where global encodings fail.
\item Exterior observers are forced to re-encode photon coherence into thermal basins.
\item Hawking radiation emerges as a basin-measure shadow of this forced re-encoding.
\end{itemize}

In this way, photon propagation, gravitational redshift, horizons, thermality, and
irreversibility are unified within a single admissibility-based framework.
% ============================================================
%  Section 8: Time, Entanglement, and Photon Encoding
% ============================================================

\section{Time, Entanglement, and Photon Encoding}
\label{sec:time_entanglement}

\subsection{Time as an encoding-dependent concept}
\label{sec:time_encoding}

Within Modal Triplet Theory, time is not a primitive parameter attached uniformly to all
degrees of freedom. Rather, time emerges as part of an effective encoding that orders
admissible updates of coherent configurations. This distinction is already implicit in the
MTT$\Rightarrow$QM and MTT$\Rightarrow$GR constructions, where the modal evolution parameter
labels updates upstairs, while physical time appears only in the effective four-dimensional
description.

In particular, time becomes meaningful when internal coherence must be rebalanced or when
records impose an ordering of events. Degrees of freedom that require no internal
rebalancing admit encodings in which no intrinsic notion of elapsed time is defined.

\subsection{Absence of intrinsic time for photons}
\label{sec:photon_time}

Photons, as massless gauge-coherence modes, carry no internal coherence inertia and therefore
no internal clock. Along admissible null updating directions, there is no accumulation of
proper time associated with the photon encoding. This is consistent with the vanishing of
proper time along null worldlines in relativistic geometry, but in MTT it is given a
structural explanation: pure gauge coherence does not require time-like rebalancing to
maintain admissibility.

Consequently, a photon encoded as a globally stretched coherence constraint does not
``experience'' time between emission and absorption. Temporal ordering appears only in the
encoding adopted by massive systems (sources, detectors, clocks) that interact with the
photon and impose time-like constraints.

\subsection{Entangled photons and shared coherence-time structure}
\label{sec:entangled_photon_time}

When photons are entangled, they are not represented as independent coherent encodings with
separate time parameters. Instead, the entangled system corresponds to a single globally
coherent configuration whose encoding does not factorize into subsystems. In this regime,
there is no meaningful notion of separate photon times or a shared clock in the usual sense.

Rather, the entangled photons share a common coherence-time structure that is effectively
atemporal until interactions or measurements enforce a partition. This explains why
entanglement correlations appear independent of spatial separation and why no time delay is
associated with the establishment of correlations: there is no internal temporal degree of
freedom to synchronize or propagate.

\subsection{Emergence of time ordering at interaction and measurement}
\label{sec:time_emergence}

Time ordering becomes unavoidable when a photon interacts strongly with matter or when a
measurement produces a stable record. In these situations, the encoding must partition the
global coherence into subsystems that maintain independent stability. This partitioning
introduces time-like rebalancing costs and therefore a local notion of elapsed time.

For example:
\begin{itemize}
\item Emission and absorption events are time-ordered relative to the clocks of the emitting
and detecting systems.
\item Scattering processes impose causal sequencing through local interactions.
\item Measurement enforces record-compatible encodings that necessarily carry time stamps.
\end{itemize}

Thus time appears not as a universal background parameter, but as a feature of encodings that
support stable, localized records.

\subsection{Compatibility with relativistic causality}
\label{sec:relativistic_time}

The absence of intrinsic photon time does not violate relativistic causality. All physical
updates propagate locally and causally in the emergent spacetime geometry, and no superluminal
signaling is possible. The atemporality of photon encoding reflects the absence of internal
time-like structure, not the absence of causal order in interactions.

Observers always describe events using their own local clocks, and these descriptions are
consistent across frames because the admissible null structure underlying photon propagation
is Lorentz invariant in the effective encoding.

\subsection{Wave--particle duality revisited}
\label{sec:wave_particle_time}

The time dependence of photon behavior further clarifies wave--particle duality. Wave-like
behavior corresponds to a globally coherent, atemporal encoding in which interference and
delocalization are admissible. Particle-like behavior corresponds to a localized, time-ordered
encoding enforced by measurement or interaction constraints.

These encodings are mutually exclusive because they impose incompatible requirements on
coherence and time ordering. The apparent duality is therefore a manifestation of switching
between admissible encodings rather than a fundamental ambiguity in the nature of the photon.

% ============================================================
%  Subsection: The Speed of Light as a Coherence Limit
% ============================================================

\subsection{The Speed of Light as a Coherence Limit}
\label{sec:speed_of_light}

The preceding sections have described photon propagation as null updating of a globally
stretched gauge-coherence encoding. We now make explicit what the speed of light $c$
represents in the MTT framework.

In Modal Triplet Theory, $c$ is not defined operationally as the speed of photons, nor
introduced axiomatically via Lorentz invariance. Instead, it arises as a structural property
of admissibility.

The coherent sector exists only if coherence updates can be stabilized by the joint
projector $\Pi$ and the fundamental contractivity condition remains satisfied. This imposes
a maximum admissible rate at which coherent constraints may be updated across the emergent
spacetime description. Updates exceeding this rate would outpace stabilization, rendering
the coherent encoding ill-defined.

This maximum admissible update rate is identified with the speed of light $c$.

Null directions in the emergent spacetime geometry correspond precisely to updates that
transport coherence without requiring internal rebalancing. Massless gauge-coherence modes
admit only such null updates, and therefore propagate at exactly $c$. Massive modes require
time-like rebalancing and consequently propagate at subluminal speeds. Spacelike updates,
corresponding to faster-than-null propagation, are inadmissible.

In this sense, the universality of $c$ reflects the universality of admissibility constraints
across the coherent sector. The role of $c$ in relativistic causality, electromagnetic
propagation, gravitational horizons, and thermality follows directly from its status as the
coherence limit of physical description.

Thus the speed of light is reinterpreted as:
\begin{quote}
\emph{the maximum rate at which coherent constraints can be updated while preserving
admissibility, stability, and boundedness of the effective encoding.}
\end{quote}

\subsection{Summary}
\label{sec:time_summary}

The MTT interpretation of photons yields a coherent account of time and entanglement:
\begin{itemize}
\item Time is an encoding-dependent construct associated with coherence rebalancing and
record formation.
\item Photons, as massless gauge-coherence modes, admit atemporal encodings along null
updating directions.
\item Entangled photons share a single global coherence structure without separate internal
times.
\item Temporal ordering emerges only when interactions or measurements enforce partition.
\end{itemize}

This perspective integrates naturally with the admissibility--encoding framework and
completes the structural account of photon behavior across quantum, relativistic, and
gravitational regimes.
% ============================================================
%  Section 9: Discussion and Outlook
% ============================================================

\section{Discussion and Outlook}
\label{sec:discussion_outlook}

\subsection{What has (and has not) been changed}
\label{sec:what_changed}

The account presented in this paper does not alter the equations of quantum electrodynamics,
general relativity, or quantum field theory on curved spacetime. Maxwell's equations, null
propagation, gravitational lensing, redshift, and Hawking's thermal spectrum remain exactly
as in the standard formulations.

What changes is the \emph{structural interpretation} of these equations within Modal Triplet
Theory. Photons are no longer treated as localized particles whose behavior must be reconciled
piecemeal across different regimes. Instead, they are understood as massless gauge-coherence
modes whose preferred encoding is global and non-factorizing, with localization, particle-like
behavior, and thermality arising only when admissibility constraints enforce a change of
encoding.

This distinction allows phenomena that are usually treated separately—wave--particle
duality, entanglement, redshift, horizons, and measurement—to be understood as manifestations
of a single underlying mechanism.

\subsection{Relation to standard interpretations}
\label{sec:relation_interpretations}

The MTT account of photons differs from common interpretations in several respects:

\begin{itemize}
\item It does not require an observer-dependent collapse postulate; encoding changes are
triggered by admissibility constraints rather than by observation.
\item It does not invoke many-worlds branching; globally coherent configurations are unique
upstairs, and apparent branching reflects non-invertibility of the observable map.
\item It does not treat thermality or information loss as fundamental; both arise from the
use of restricted encodings across admissibility barriers.
\end{itemize}

At the same time, the MTT picture is compatible with algebraic QFT insights regarding
non-factorization of local algebras and with semiclassical derivations of Hawking radiation.
It should therefore be viewed as a synthesis rather than a replacement of existing results.

\subsection{Photons as a unifying case study}
\label{sec:photons_unifying}

Photons provide a particularly transparent case study because they sit at the intersection
of quantum theory, relativity, and gravity. Their lack of internal inertia exposes the
distinction between global coherence and local encoding more clearly than in massive systems.

In the MTT framework, photons illustrate how:
\begin{itemize}
\item global entanglement can be the preferred physical encoding,
\item null propagation arises from admissibility rather than kinematic assumption,
\item localization is conditional and measurement-driven,
\item horizons act as encoding barriers rather than physical singularities.
\end{itemize}

These lessons generalize to other gauge fields and, with appropriate modifications, to
gravitational waves and massive particles.

\subsection{Implications for quantum gravity and information}
\label{sec:implications_qg}

By treating horizons as admissibility barriers and Hawking radiation as forced re-encoding,
MTT offers a structural resolution of the black hole information problem. Information is not
destroyed at the modal level; it becomes inaccessible to a particular encoding. This places
black hole thermality and measurement-induced irreversibility on the same conceptual footing.

More broadly, the admissibility--encoding framework suggests that many puzzles traditionally
framed as paradoxes arise from insisting on a single encoding outside its domain of validity.
Recognizing and respecting encoding boundaries may therefore be essential for progress in
quantum gravity.

\subsection{Limitations and open questions}
\label{sec:limitations}

The present work is primarily interpretive and structural. Several important questions remain
open:
\begin{itemize}
\item A fully explicit microphysical derivation of horizon parameters (e.g.\ the Hawking
temperature) within MTT remains to be developed.
\item The quantitative limits of coherence persistence and selection-front behavior in
realistic environments require further analysis.
\item Extensions to interacting gauge theories and strongly coupled regimes merit separate
treatment.
\end{itemize}

Addressing these questions will require combining the encoding perspective developed here
with detailed calculations in specific admissible realizations.

\subsection{Outlook}
\label{sec:outlook}

The interpretation of photons presented in this paper illustrates how Modal Triplet Theory
can unify diverse physical phenomena through a single admissibility-based framework. By
emphasizing encoding over ontology and projection over postulate, MTT provides a language in
which quantum mechanics, relativity, and thermality appear as complementary descriptions of
the same underlying coherent structure.

Future work will extend this analysis to other fields and interactions, and will explore how
the admissibility--encoding viewpoint can guide both conceptual understanding and practical
calculations in quantum gravity, cosmology, and quantum information.

\begin{thebibliography}{99}

\bibitem{MTT-Admissibility}
P.~Nero,
\emph{Modal Triplet Theory: Admissibility, Encodings, and the Structure of Physical Description},
Zenodo preprint, January 2026.
\url{https://doi.org/10.5281/zenodo.18255621}

\bibitem{MTT-Foundation}
P.~Nero,
\emph{Modal Triplet Theory: Foundation},
Zenodo preprint, September 2025.
\url{https://doi.org/10.5281/zenodo.16949762}

\bibitem{MTT-FixedPoints}
P.~Nero,
\emph{Fixed Points I--VI: Complete Coherence Spine},
Zenodo preprints, August 2025.
\url{https://doi.org/10.5281/zenodo.16948748}

\bibitem{MTT-ProjectionAdmissibility}
P.~Nero,
\emph{The Projection--Admissibility Principle: Structural Constraints on Effective Physical Description},
Zenodo preprint, January 2026.
\url{https://doi.org/10.5281/zenodo.18255838}

\bibitem{MTT-Closure}
P.~Nero,
\emph{Closure and Inevitability in Modal Triplet Theory},
Zenodo preprint, January 2026.
\url{https://doi.org/10.5281/zenodo.18255510}

\bibitem{MTT-CoherenceCapacity}
P.~Nero,
\emph{Coherence Capacity as the Fundamental Resource of Effective Physics},
Zenodo preprint, January 2026.
\url{https://doi.org/10.5281/zenodo.18255905}

\bibitem{MTT-CoherenceDynamics}
P.~Nero,
\emph{Dynamics of Coherence Capacity: Transport, Concentration, and Exhaustion},
Zenodo preprint, January 2026.
\url{https://doi.org/10.5281/zenodo.18256048}

\bibitem{MTT-QM}
P.~Nero,
\emph{Modal Triplet Theory: From MTT to Quantum Mechanics},
Zenodo preprint, September 2025.
\url{https://doi.org/10.5281/zenodo.17074246}

\bibitem{MTT-QFT}
P.~Nero,
\emph{From MTT to Quantum Field Theory},
Zenodo preprint, 2025.
\url{https://doi.org/10.5281/zenodo.17068816}

\bibitem{MTT-GR}
P.~Nero,
\emph{Modal Triplet Theory: From MTT to General Relativity},
Zenodo preprint, October 2025.
\url{https://doi.org/10.5281/zenodo.16950597}

\bibitem{MTT-UVQG}
P.~Nero,
\emph{Modal Triplet Theory: From MTT to a UV-Finite, Unitary Quantum Gravity},
Zenodo preprint, 2025.
\url{https://doi.org/10.5281/zenodo.17077671}

\bibitem{MTT-Measurement}
P.~Nero,
\emph{Measurement as Disturbance and Stabilization in Modal Triplet Theory},
Zenodo preprint, 2025.
\url{https://doi.org/10.5281/zenodo.17177404}

\bibitem{MTT-Irreversibility}
P.~Nero,
\emph{Projection, Probability, and Irreversibility: Shadow Bridges Between Measurement, Black Holes, and Cosmology in Modal Triplet Theory},
Zenodo preprint, January 2026.
\url{https://doi.org/10.5281/zenodo.18256408}

\bibitem{MTT-Bell-Beables}
P.~Nero,
\emph{Modal Fixed Points, Bell’s Beables, and the Limits of Factorization},
Zenodo preprint, 2025.
\url{https://doi.org/10.5281/zenodo.17076300}

\bibitem{MTT-Bell-Temporal}
P.~Nero,
\emph{Temporal Bell Inequalities and Global Consistency in Modal Triplet Theory},
Zenodo preprint, August 2025.
\url{https://doi.org/10.5281/zenodo.18208884}

\bibitem{MTT-Stochastic}
P.~Nero,
\emph{From Modal Triplet Theory to Indivisible Stochastic Processes: A First-Principles, Fully Rigorous Derivation},
Zenodo preprint, January 2026.
\url{https://doi.org/10.5281/zenodo.18254862}

\end{thebibliography}
\end{document}
